\documentclass[10pt,a4paper]{amsart}
\usepackage[margin=19mm]{geometry}
\usepackage{amsmath,amssymb,mathtools}
\usepackage[scr=rsfs,cal=euler]{mathalfa}
\usepackage{xfrac}
\usepackage[unicode,hidelinks,bookmarks=false]{hyperref}
\newcommand{\ie}{i.e.\ }
\newcommand{\cf}{cf.\ }
\newcommand{\resp}{respectively\ }
\newcommand{\Subsection}{\subsection}
\providecommand{\nobreakdash}{\nobreak\hbox{-}}
\setlength{\emergencystretch}{2em}
\multlinegap0pt
\usepackage{float}
\usepackage{amssymb}
\usepackage{xcolor}
\newtheorem{remark}{Remark}
\newtheorem{lemma}{Lemma}
\title{Chaotic Dynamics of Conformable Semigroups via Classical Theory}
\author{Mohamed Khoulane, Aziz El Ghazouani, M'hamed Elomari}
\date{Source: arXiv:2602.06782v1 (CC BY 4.0)}
\begin{document}
\maketitle

\noindent\emph{Source and licence.} Chaotic Dynamics of Conformable Semigroups via Classical Theory. Version arXiv:2602.06782v1 (CC BY 4.0), distributed by arXiv under the Creative Commons Attribution 4.0 International licence, http://creativecommons.org/licenses/by/4.0/, as recorded in the arXiv licence metadata for this exact version and shown on the abstract page https://arxiv.org/abs/2602.06782v1. Full licence notice: \texttt{licenses/arxiv-2602.06782-CC-BY-4.0.txt}.\par\smallskip

\noindent\textit{Editorial scope.} Counted unit: the article's lemma lem:invariance\_sets (source0.tex lines 956--963). The article's complete original proof of it is delivered with it (source0.tex lines 965--1019, one \texttt{proof} environment of 55 lines). No mathematics is written, repaired or re-derived: the statement and the proof are the article's own lines, in the article's own notation, and no retained line is substituted or re-flowed.  Retained uncounted as the source-local context the proof needs: lines 832--839; lines 841--843; lines 889--899.  Nothing was omitted from any retained span, so the package carries no omission record.  No retained line contains a citation, so the wrapper carries no bibliography and no bibliography entry.  No retained line contains a figure environment, an image-inclusion command or drawing code, so no image is delivered and the package declares no asset dependency.  Editorial formatting only, in addition: an AMS \texttt{amsart} preamble that restates the article's own declarations and macros used by the retained lines, namely the environments \texttt{remark}, \texttt{lemma} on separate counters, together with the standard packages those lines need.  The retained environments print in this document as Remark 1, Lemma 1 in order of appearance, whereas the article prints the same items with the numbering of its own full document.  The complete native source of the article as distributed in the arXiv e-print for this exact version is bundled unchanged as \texttt{sources/194121/source0.tex}.
\begin{equation}\label{eq:defT}
\mathcal T(s)
:=
\mathcal S_\delta\!\bigl(\Psi^{-1}(s)\bigr)
=
\mathcal S_\delta\!\bigl((\delta s)^{1/\delta}\bigr),
\qquad s\ge0,
\end{equation}

\begin{equation}\label{eq:conjugacy_time}
\mathcal S_\delta(t)=\mathcal T(\Psi(t)),\qquad t\ge0.
\end{equation}

\begin{remark}[Orbit identity]\label{rem:orbit_identity}
From \eqref{eq:conjugacy_time} and the fact that $\Psi$ is a bijection of $[0,\infty)$,
we infer that for each $x\in X$,
\[
\{\mathcal S_\delta(t)x:\ t\ge0\}
=
\{\mathcal T(s)x:\ s\ge0\}.
\]
Hence conformable reparametrization changes only the parametrization of the orbit,
not the set of points visited along it.
\end{remark}

\begin{lemma}[Invariance of dynamical sets]\label{lem:invariance_sets}
With $\mathcal T$ defined by \eqref{eq:defT}, one has
\[
X_0^\delta=X_0,\qquad X_\infty^\delta=X_\infty,\qquad X_p^\delta=X_p.
\]
In particular, $\mathcal S_\delta$ is $\delta$--hypercyclic (resp.\ $\delta$--chaotic)
if and only if $\mathcal T$ is hypercyclic (resp.\ chaotic) in the classical sense.
\end{lemma}

\begin{proof}
We use \eqref{eq:conjugacy_time} and the fact that $\Psi$ is a homeomorphism of
$[0,\infty)$ onto itself.

\smallskip
\noindent\emph{Step 1: invariance of $X_0^\delta$.}
For $x\in X$,
\[
\lim_{t\to\infty}\mathcal S_\delta(t)x=0
\ \Longleftrightarrow\
\lim_{t\to\infty}\mathcal T(\Psi(t))x=0.
\]
Since $\Psi(t)\to\infty$ as $t\to\infty$ and $\Psi([0,\infty))=[0,\infty)$,
this is equivalent to $\lim_{s\to\infty}\mathcal T(s)x=0$.
Hence $X_0^\delta=X_0$.

\smallskip
\noindent\emph{Step 2: invariance of periodic points.}
If $x\in X_p^\delta$, then $\mathcal S_\delta(t)x=x$ for some $t>0$, and thus
$\mathcal T(\Psi(t))x=x$ with $\Psi(t)>0$, so $x\in X_p$.
Conversely, if $x\in X_p$, there exists $s>0$ with $\mathcal T(s)x=x$; choosing
$t=\Psi^{-1}(s)$ yields
\[
\mathcal S_\delta(t)x=\mathcal T(\Psi(t))x=\mathcal T(s)x=x,
\]
so $x\in X_p^\delta$.
Therefore $X_p^\delta=X_p$.

\smallskip
\noindent\emph{Step 3: invariance of $X_\infty^\delta$.}
Fix $x\in X$ and $\varepsilon>0$.
Assume $x\in X_\infty^\delta$. Then there exist $y\in X$ and $t>0$ such that
$\|y\|<\varepsilon$ and $\|\mathcal S_\delta(t)y-x\|<\varepsilon$.
Since $\mathcal S_\delta(t)=\mathcal T(\Psi(t))$, we get
$\|\mathcal T(\Psi(t))y-x\|<\varepsilon$, hence $x\in X_\infty$.
Conversely, if $x\in X_\infty$, there exist $y\in X$ and $s>0$ such that
$\|y\|<\varepsilon$ and $\|\mathcal T(s)y-x\|<\varepsilon$.
Setting $t=\Psi^{-1}(s)$ yields
\[
\|\mathcal S_\delta(t)y-x\|
=
\|\mathcal T(\Psi(t))y-x\|
=
\|\mathcal T(s)y-x\|
<\varepsilon,
\]
so $x\in X_\infty^\delta$.
Thus $X_\infty^\delta=X_\infty$.

\smallskip
\noindent\emph{Step 4: hypercyclicity and chaos.}
By Remark~\ref{rem:orbit_identity}, orbit sets coincide for $\mathcal S_\delta$ and
$\mathcal T$, hence the existence of a dense orbit is equivalent for the two families.
Together with $X_p^\delta=X_p$, this yields the equivalence of chaos.
\end{proof}

\end{document}
